\documentclass{article}

\usepackage[submission]{../../../colm_final/paper/colm2026_conference}
\usepackage{booktabs}
\usepackage{hyperref}
\usepackage{url}

\title{ThoughtFlow: A Sparse-Cache Falsification Study for Reasoning-Trace Retention}
\author{Anonymous authors\\Paper under double-blind review}

\begin{document}
\maketitle

\begin{abstract}
The ThoughtFlow-FP8 project tests whether training-free sparse KV retention can preserve useful reasoning context better than matched-budget KV-retention proxies. The current contribution is diagnostic, not a positive method, and no real FP8, CUDA, latency, or throughput result is claimed. Synthetic traces support the intuition that phase markers and anchors can be protected, but real retained-context quality and sparse-cache gates rule out further tuning of the original anchor/recent/phase/math family on the available traces. We then pre-registered three successor signals. Recurrence-distance utility (\texttt{rdu\_topk}) initially reaches NLL 3.779 versus ThinKV-like 3.900 and R-KV-like 3.939 on a frozen 74-trace distilgpt2 gate, but reproduction gates demote it. Prefix-surprisal utility (\texttt{psi\_topk}) fails on a fresh 70-trace C2C GSM surface, reaching NLL 7.899 versus ThinKV-like 3.906 and R-KV-like 3.960. Value-weighted attention contribution (\texttt{vwac\_topk}) fails on a fresh 64-trace C2C SVAMP surface, reaching NLL 4.336 versus R-KV-like 4.096 and ThinKV-like 4.162. The anchor/phase int8 quantization primitive passes Triton interpreter correctness tests against its CPU reference. The paper therefore reports a sparse-cache falsification ladder and no live method branch.
\end{abstract}

\section{Policy Family}

The original policy protects anchors and phase labels. The first successor adds a small recent-token reserve. The second successor protects anchors, reserves half the retained budget for recent tokens, and lets phase markers compete with math-state and high-importance reasoning tokens under a saliency bonus. Those policies are now stopped on the current saved traces. The pre-registered \texttt{rdu\_topk} successor is demoted after stricter reproduction; later \texttt{psi\_topk} and \texttt{vwac\_topk} successors are killed by fresh-surface gates.

\begin{table}[h]
\centering
\small
\begin{tabular}{lrrrr}
\toprule
Policy & Keep rate & NLL & $\Delta$ NLL vs full & PPL \\
\midrule
full context & 1.000 & 2.101 & 0.000 & 9.7 \\
R-KV-like & 0.210 & 3.419 & 1.319 & 38.7 \\
ThinKV-like & 0.210 & 3.583 & 1.482 & 44.6 \\
LongFlow-like & 0.210 & 3.961 & 1.861 & 74.2 \\
ThoughtFlow & 0.210 & 3.961 & 1.861 & 74.2 \\
ThoughtFlow-recent & 0.210 & 3.562 & 1.461 & 44.2 \\
ThoughtFlow-saliency-recent & 0.210 & 3.434 & 1.333 & 38.1 \\
Held-out R-KV-like & 0.211 & 3.482 & -- & 43.7 \\
Held-out sweep best & 0.211 & 3.480 & -- & 43.7 \\
\bottomrule
\end{tabular}
\caption{Retained-context continuation proxy on distilgpt2 traces. Lower NLL is better. Full context is not a matched-budget baseline. The held-out sweep row is selected on 12 train traces and reported on 12 held-out traces, so its NLL is not directly averaged with the all-trace rows.}
\end{table}

\section{Hidden/KV Telemetry}

We then ran a Mac-local telemetry gate using real distilgpt2 tensors rather than only text markers. The probe ranks retained tokens by attention received, final-hidden norm, key norm, value norm, or combined KV norm. This is still not sparse-KV decoding, but it checks whether the interpretable phase policy is merely preserving labels that hidden/KV saliency would not select.

\begin{table}[h]
\centering
\small
\begin{tabular}{lrrrr}
\toprule
Policy & Keep rate & Anchor & Phase & Math-state \\
\midrule
ThoughtFlow & 0.207 & 1.000 & 1.000 & 0.918 \\
LongFlow-like & 0.207 & 1.000 & 1.000 & 0.918 \\
ThinKV-like & 0.207 & 1.000 & 0.948 & 0.848 \\
\texttt{value\_norm\_topk} & 0.207 & 0.000 & 0.492 & 0.845 \\
train-selected sweep & 0.207 & 1.000 & 0.430 & 0.956 \\
\bottomrule
\end{tabular}
\caption{Hidden/KV saliency telemetry on 24 distilgpt2 traces. The strongest real-saliency proxy by phase+math recall is \texttt{value\_norm\_topk}. ThoughtFlow beats it on phase recall, but LongFlow-like ties ThoughtFlow under the local importance heuristic.}
\end{table}

Paired against \texttt{value\_norm\_topk}, ThoughtFlow's phase-recall delta is +0.508 with 95\% CI [+0.436, +0.579]. The math-state delta is only +0.073 with 95\% CI [-0.078, +0.223]. This rules out a positive claim based on phase recall alone: it may be marker preservation rather than useful cache retention.

\section{CPU Sparse-Cache Quality Gate}

Finally, we ran a CPU-only sparse-cache probe. The model processes the full
prefix once, each policy prunes the returned KV cache to the same 0.20 budget,
and the continuation is scored from the pruned cache. This is not a Triton,
CUDA, FP8, latency, or throughput result, but it is stronger than retained-text
NLL because the model actually receives a sparse cache.

\begin{table}[h]
\centering
\small
\begin{tabular}{lrrr}
\toprule
Policy & Keep rate & NLL & Paired $\Delta$ vs R-KV-like \\
\midrule
full cache & 1.000 & 2.142 & -1.296 [-1.533,-1.066] \\
ThoughtFlow-saliency-recent & 0.210 & 3.372 & -0.067 [-0.151,+0.011] \\
ThinKV-like & 0.210 & 3.389 & -0.049 [-0.192,+0.077] \\
ThoughtFlow-recent & 0.210 & 3.399 & -0.040 [-0.169,+0.074] \\
ThoughtFlow sweep best & 0.210 & 3.436 & -0.002 [-0.007,+0.000] \\
R-KV-like & 0.210 & 3.438 & 0.000 \\
LongFlow-like & 0.210 & 3.588 & +0.150 [-0.016,+0.300] \\
ThoughtFlow & 0.210 & 3.588 & +0.150 [-0.006,+0.297] \\
\bottomrule
\end{tabular}
\caption{CPU sparse-KV continuation quality on 24 distilgpt2 traces. Negative paired delta is better than R-KV-like. ThoughtFlow-saliency-recent is best in mean NLL, but it clears the strongest non-ThoughtFlow row by only 0.017 NLL and remains below the pre-registered 0.03 promotion margin.}
\end{table}

A bounded sparse-cache sweep over 24 nearby anchor/recent/phase/math
configurations selects \texttt{tf\_sparse\_r0.55\_p0.05\_m0.12\_a2} on the
12 train traces. On the 12 held-out traces it gets NLL 3.340 versus ThinKV-like
3.385 and R-KV-like 3.420. This clears the mean 0.03 margin, and the paired
delta versus R-KV-like is -0.080 with 95\% CI [-0.152,-0.014]. However, the
paired delta versus ThinKV-like is -0.045 with 95\% CI [-0.226,+0.182]. The
fixed ThoughtFlow-saliency-recent incumbent has lower held-out mean NLL, 3.304,
but its paired CIs also cross zero. This is a promising falsification candidate,
not a positive result.

We then froze \texttt{ThoughtFlow-saliency-recent} and
\texttt{tf\_sparse\_r0.55\_p0.05\_m0.12\_a2} and ran a larger 74-trace sparse
cache slice with no retuning. ThinKV-like is the best compressed row among the
stopped family comparison at NLL 3.900. The frozen sparse ThoughtFlow row gets
3.908, with paired delta -0.031 versus R-KV-like and 95\% CI
[-0.078,+0.020], but +0.008 versus ThinKV-like with 95\% CI
[-0.060,+0.085]. ThoughtFlow-saliency-recent gets 3.920. The larger
no-retuning gate therefore stops this policy family.

After that stop decision, we pre-registered \texttt{rdu\_topk}. It computes
recurrence-distance utility from full-prefix prefill attentions using lag buckets
[8,15], [16,31], [32,63], and [64,$\infty$], then keeps the top budget by utility
with ties broken by lower position. It does not use continuation tokens, trace
labels, recent reserves, anchors, phase markers, or math-state bonuses. After
\texttt{rdu\_topk} failed reproduction, we pre-registered \texttt{psi\_topk}
(high self-surprisal prefix tokens) and \texttt{vwac\_topk} (future attention
weighted by cached value-vector norm).

\section{Evidence Ladder}

\begin{table}[h]
\centering
\small
\begin{tabular}{p{0.28\linewidth}p{0.42\linewidth}p{0.20\linewidth}}
\toprule
Gate & Readout & Decision \\
\midrule
Synthetic retention & Anchor/phase labels can be protected at matched keep rate. & Intuition only. \\
Retained-text NLL & Original ThoughtFlow rows do not beat the strongest proxy. & Weakened. \\
Hidden/KV telemetry & Phase recall improves, but math-state impact is not stable and LongFlow-like ties the heuristic. & Diagnostic. \\
CPU sparse-cache probe & A tuned sparse ThoughtFlow row is promising in mean but not robust against ThinKV-like uncertainty. & Not promoted. \\
Frozen 74-trace gate & \texttt{rdu\_topk} beats ThinKV-like and R-KV-like on the first frozen surface. & Candidate only. \\
Same-slice rerun & Cached result reproduces exactly on the deterministic Mac stack. & Repro bookkeeping. \\
Alternate surface & Cross-family rows still lose, but a stopped same-family row beats \texttt{rdu\_topk} by 0.006 NLL. & Weakened. \\
Independent saved traces & R-KV-like is best compressed at NLL 3.981; \texttt{rdu\_topk} reaches 4.014. & Killed. \\
Prefix-surprisal fresh surface & \texttt{psi\_topk} reaches NLL 7.899 versus ThinKV-like 3.906 and R-KV-like 3.960. & Killed. \\
Value-weighted attention fresh surface & \texttt{vwac\_topk} reaches NLL 4.336 versus R-KV-like 4.096 and ThinKV-like 4.162. & Killed. \\
\bottomrule
\end{tabular}
\caption{ThoughtFlow-FP8 evidence ladder. The positive-looking gate is shown as one step that failed to reproduce, not as the headline result.}
\end{table}

\begin{table}[h]
\centering
\small
\begin{tabular}{lrrrr}
\toprule
Policy & Traces & Keep rate & NLL & Paired $\Delta$ vs ThinKV-like \\
\midrule
full cache & 74 & 1.000 & 2.848 & -1.052 [-1.199,-0.919] \\
\texttt{rdu\_topk} & 74 & 0.213 & 3.779 & -0.121 [-0.211,-0.037] \\
ThinKV-like & 74 & 0.213 & 3.900 & 0.000 \\
frozen sparse ThoughtFlow & 74 & 0.213 & 3.908 & +0.008 [-0.060,+0.085] \\
ThoughtFlow-saliency-recent & 74 & 0.213 & 3.920 & +0.020 [-0.030,+0.074] \\
R-KV-like & 74 & 0.214 & 3.939 & +0.039 [-0.015,+0.099] \\
LongFlow-like & 74 & 0.213 & 4.158 & +0.258 [+0.178,+0.337] \\
\bottomrule
\end{tabular}
\caption{Historical one-shot frozen sparse-cache successor evaluation. \texttt{rdu\_topk} also beats R-KV-like by -0.160 NLL with 95\% CI [-0.264,-0.050] on this first surface, but later alternate-surface and independent-trace gates fail to reproduce the promotion.}
\end{table}

\paragraph{Cached and same-slice checks.}
Even/odd and half-split diagnostics on the first 74-trace surface preserve the
\texttt{rdu\_topk} mean margin, but only 2/4 half-size partitions keep both
paired CIs below zero. A measured same-slice rerun exactly reproduces the cached
NLLs on the deterministic Mac stack and adds oracle/headroom diagnostics
(\texttt{rdu\_topk} is 0.145 NLL above the compressed oracle), but it is not an
independent reproduction.

\paragraph{Measured alternate-surface check.}
We then changed only the measurement surface, using \texttt{max\_length=112}
and \texttt{continuation\_tokens=32}, while preserving the same
\texttt{rdu\_topk} rule, matched 0.20 keep fraction, cached-versus-measured
labels, same-family/cross-family reporting, and oracle/headroom diagnostics.
This check is not a clean reproduction. \texttt{rdu\_topk} remains strong
against cross-family rows, reaching NLL 3.594 versus R-KV-like 3.681 and
ThinKV-like 3.851, with paired deltas -0.087 [-0.139,-0.028] and -0.256
[-0.465,-0.086]. However, \texttt{tf\_sparse\_r0.55\_p0.05\_m0.12\_a2}, a
stopped same-family sparse row, reaches NLL 3.588 and beats \texttt{rdu\_topk}
by 0.006. The measured per-trace compressed oracle reaches NLL 3.460, leaving
\texttt{rdu\_topk} 0.135 above oracle and 0.847 above full cache, with a 0.438
oracle-hit rate.

\paragraph{Independent saved-trace reproduction check.}
The stronger no-retuning check changes the saved trace inputs to a 96-row chat
SVAMP surface while keeping the same \texttt{rdu\_topk} rule and 0.20 keep
fraction. It scores 89 traces after token-length filtering. This does not
reproduce the positive gate. R-KV-like is now the best compressed row at NLL
3.981, while \texttt{rdu\_topk} reaches 4.014. The paired delta versus
R-KV-like is +0.032 with 95\% CI [-0.071,+0.137], and ThinKV-like is only
0.030 NLL worse than \texttt{rdu\_topk} with a confidence interval crossing
zero. Same-family stopped rows are worse than \texttt{rdu\_topk}, but only by
0.006 and 0.010 NLL. The compressed per-trace oracle reaches NLL 3.754, leaving
\texttt{rdu\_topk} 0.260 above oracle and 1.083 above full cache. This rules
out claiming \texttt{rdu\_topk} as a robust positive method on the available
surfaces.

A coarse trace-level decomposition explains where the reproduction fails:
\texttt{rdu\_topk} is +0.213 NLL worse than R-KV-like on long-prefix/high-RDU-density
rows but -0.049 better on short-prefix/low-density rows. The large-regression
buckets are exactly the buckets a robust policy would need to handle before any
long-context claim.

\paragraph{Prefix-surprisal successor.}
The final Mac-feasible successor tested whether retaining high self-surprisal
prefix tokens could preserve information that is hard to reconstruct. The signal
uses only prefill logits, not continuation loss, trace labels, recent reserves,
or RDU lag buckets. On a fresh 70-trace C2C GSM saved-trace surface,
\texttt{psi\_topk} fails decisively: ThinKV-like is best compressed at NLL 3.906,
R-KV-like reaches 3.960, and \texttt{psi\_topk} reaches 7.899. Paired deltas are
+3.939 [+3.720,+4.144] versus R-KV-like and +3.994 [+3.773,+4.205] versus
ThinKV-like, where positive is worse for \texttt{psi\_topk}. This rules out this
exact signal in the current Mac sparse-cache harness.

\paragraph{Value-weighted attention successor.}
The final successor tested a cache-state utility: future prefix attention to a
token multiplied by the token's cached value-vector norm. On a fresh 64-trace
C2C SVAMP surface, \texttt{vwac\_topk} also fails. R-KV-like is best compressed at
NLL 4.096, ThinKV-like reaches 4.162, and \texttt{vwac\_topk} reaches 4.336.
Paired deltas are +0.241 [+0.030,+0.530] versus R-KV-like and +0.174
[+0.002,+0.392] versus ThinKV-like.

\section{Decision}

The branch is now demoted to a diagnostic. The stopped anchor/recent/phase/math
policy family remains ruled out as a tunable branch on this trace set,
\texttt{rdu\_topk} fails reproduction, and \texttt{psi\_topk}/\texttt{vwac\_topk}
fail their fresh-surface gates. A paper-ready positive claim now requires a genuinely new
pre-registered utility signal evaluated once on a fresh/larger frozen surface
with strict same-family, cross-family, and oracle/headroom reporting. The Triton
interpreter gate is no longer the blocker for the quantization scaffold; method
evidence is.

\section{Related Work and Scope}

This draft intentionally does not claim a new state-of-the-art KV-compression method. ThinKV, LongFlow, R-KV/R-KVHash-style retention, LazyEviction/ForesightKV-style future-use ideas, PM-KVQ, and KVQuant already occupy much of the sparse-cache and quantized-cache design space. Our contribution is narrower: a Mac-local falsification ladder that shows how plausible interpretable retention signals can clear one frozen surface, fail stricter reproduction, or collapse on fresh surfaces. The project name retains ``FP8'' because that was the original systems target, but the current artifacts only validate CPU sparse-cache scoring and an int8/Triton-interpreter reference primitive; they do not measure FP8 numerical drift or native FP8 serving.

\section{Reproducibility}

All reported gates use tracked scripts under \url{experimental/thoughtflow_fp8/phase2/} and the repo-local \texttt{./venv\_arm64}. The frozen sparse-cache probe, same-slice rerun, alternate-surface rerun, independent saved-trace rerun, and fresh-surface successor checks each write Markdown and JSON artifacts with promotion criteria, paired intervals, oracle/headroom diagnostics, and no-retuning decisions. The stable local correctness command is recorded in the reviewer pack.

\section{Artifacts}

Key artifacts are the current decision manifest, reviewer pack, RDU independent
reproduction check, PSI/VWAC fresh-surface checks, and Triton install note:
\url{experimental/thoughtflow_fp8/phase2/current_decision_manifest_20260506.md},
\url{experimental/thoughtflow_fp8/paper/reviewer_pack.md},
\url{experimental/thoughtflow_fp8/phase2/rdu_independent_trace_reproduction_check.md},
\url{experimental/thoughtflow_fp8/phase2/psi_fresh_sparse_cache_check.md},
\url{experimental/thoughtflow_fp8/phase2/vwac_fresh_sparse_cache_check.md}, and
\url{experimental/triton_cpu_source_install_20260506.md}.

\end{document}
